\documentclass{article}

\usepackage{fancyhdr}
\usepackage{extramarks}
\usepackage{amsmath}
\usepackage{amsthm}
\usepackage{amsfonts}
\usepackage{tikz}
\usepackage[plain]{algorithm}
\usepackage{algpseudocode}
\usepackage{listings}
\usepackage{xcolor}

\definecolor{codegreen}{rgb}{0,0.6,0}
\definecolor{codegray}{rgb}{0.5,0.5,0.5}
\definecolor{codepurple}{rgb}{0.58,0,0.82}
\definecolor{backcolour}{rgb}{0.95,0.95,0.92}

\lstset{
    backgroundcolor=\color{backcolour},
    commentstyle=\color{codegreen},
    keywordstyle=\color{magenta},
    numberstyle=\tiny\color{codegray},
    stringstyle=\color{codepurple},
    basicstyle=\ttfamily\footnotesize,
    breakatwhitespace=false,
    breaklines=true,
    captionpos=b,
    keepspaces=true,
    numbers=left,
    numbersep=5pt,
    showspaces=false,
    showstringspaces=false,
    showtabs=false,
    tabsize=2
}

\usetikzlibrary{automata,positioning}

%
% Basic Document Settings
%

\topmargin=-0.45in
\evensidemargin=0in
\oddsidemargin=0in
\textwidth=6.5in
\textheight=9.0in
\headsep=0.25in

\linespread{1.1}

\pagestyle{fancy}
\lhead{\hmwkAuthorName}
\chead{\hmwkClass\ (\hmwkClassInstructor\ \hmwkClassTime): \hmwkTitle}
\rhead{\firstxmark}
\lfoot{\lastxmark}
\cfoot{\thepage}

\renewcommand\headrulewidth{0.4pt}
\renewcommand\footrulewidth{0.4pt}

\setlength\parindent{0pt}

%
% Create Problem Sections
%

\newcommand{\enterProblemHeader}[1]{
	\nobreak\extramarks{}{Problem \arabic{#1} continued on next page\ldots}\nobreak{}
	\nobreak\extramarks{Problem \arabic{#1} (continued)}{Problem \arabic{#1} continued on next page\ldots}\nobreak{}
}

\newcommand{\exitProblemHeader}[1]{
	\nobreak\extramarks{Problem \arabic{#1} (continued)}{Problem \arabic{#1} continued on next page\ldots}\nobreak{}
	\stepcounter{#1}
	\nobreak\extramarks{Problem \arabic{#1}}{}\nobreak{}
}

\setcounter{secnumdepth}{0}
\newcounter{partCounter}
\newcounter{homeworkProblemCounter}
\setcounter{homeworkProblemCounter}{1}
\nobreak\extramarks{Problem \arabic{homeworkProblemCounter}}{}\nobreak{}

%
% Homework Problem Environment
%
% This environment takes an optional argument. When given, it will adjust the
% problem counter. This is useful for when the problems given for your
% assignment aren't sequential. See the last 3 problems of this template for an
% example.
%
\newenvironment{homeworkProblem}[1][-1]{
	\ifnum#1>0
	\setcounter{homeworkProblemCounter}{#1}
	\fi
	\section{Problem \arabic{homeworkProblemCounter}}
	\setcounter{partCounter}{1}
	\enterProblemHeader{homeworkProblemCounter}
}{
	\exitProblemHeader{homeworkProblemCounter}
}

%
% Homework Details
%   - Title
%   - Due date
%   - Class
%   - Section/Time
%   - Instructor
%   - Author
%

\newcommand{\hmwkTitle}{Labwork 1}
\newcommand{\hmwkDueDate}{DUE DATE}
\newcommand{\hmwkClass}{Deep Learning}
\newcommand{\hmwkClassTime}{}
\newcommand{\hmwkClassInstructor}{SonTG}
\newcommand{\hmwkAuthorName}{\textbf{Nguyen Trong Minh}}

%
% Title Page
%

\title{
	\vspace{2in}
	\textmd{\textbf{\hmwkClass:\ \hmwkTitle}}\\
	%\normalsize\vspace{0.1in}\small{Due\ on\ \hmwkDueDate}\\
	\vspace{0.1in}\large{\textit{\hmwkClassInstructor\ \hmwkClassTime}}
	\vspace{3in}
}

\author{\hmwkAuthorName}
\date{}

\renewcommand{\part}[1]{\textbf{\large Part \Alph{partCounter}}\stepcounter{partCounter}\\}

%
% Various Helper Commands
%

% Useful for algorithms
\newcommand{\alg}[1]{\textsc{\bfseries \footnotesize #1}}

% For derivatives
\newcommand{\deriv}[1]{\frac{\mathrm{d}}{\mathrm{d}x} (#1)}

% For partial derivatives
\newcommand{\pderiv}[2]{\frac{\partial}{\partial #1} (#2)}

% Integral dx
\newcommand{\dx}{\mathrm{d}x}

% Alias for the Solution section header
\newcommand{\solution}{\textbf{\large Solution}}

% Probability commands: Expectation, Variance, Covariance, Bias
\newcommand{\E}{\mathrm{E}}
\newcommand{\Var}{\mathrm{Var}}
\newcommand{\Cov}{\mathrm{Cov}}
\newcommand{\Bias}{\mathrm{Bias}}

\begin{document}
	
	\maketitle
	\pagebreak
	
	\begin{homeworkProblem}
		\solution
		
		\subsection{Algorithm Implementation}
		
		The Gradient Descent algorithm is implemented to find the minimum of a quadratic function $f(x) = x^2$. The implementation consists of the following components:
		
		\begin{itemize}
			\item \textbf{Forward Function}: Computes the objective function value $f(x) = x^2$.
			\item \textbf{Gradient Function}: Computes the derivative $f'(x) = 2x$.
			\item \textbf{Update Rule}: The parameter is updated iteratively using the formula:
			\[ x_{t+1} = x_t - \eta \cdot \nabla f(x_t) \]
			where $\eta$ is the learning rate.
			\item \textbf{Stopping Criterion}: The algorithm stops when the absolute change in the function value between successive iterations falls below the learning rate $\eta$:
			\[ |f(x_{t+1}) - f(x_t)| < \eta \]
		\end{itemize}
		
		The Python code for this implementation is provided below:
		
		\begin{lstlisting}[language=Python, caption=Gradient Descent Implementation]
def forward(x):
    return x*x

def gradient(x):
    return 2*x

def descent(forward, gradient, init, lr=10**(-2)):
    value = [1,0]
    time = 0
    while ((forward(value[1]) - forward(value[0]))**2)**(1/2) >= lr:
        print(time, ": ", init, " ", forward(init)) 
        time +=1
        init -= gradient(init) * lr
        value[0] = value[1]
        value[1] = init

descent(forward, gradient, init = 10)
		\end{lstlisting}
		
		\subsection{Effect of Learning Rate Analysis}
		
		The learning rate ($\eta$) is a critical hyperparameter that determines the step size at each iteration. For the function $f(x) = x^2$, the update equation can be simplified as:
		\[ x_{t+1} = x_t - \eta(2x_t) = x_t(1 - 2\eta) \]
		
		The behavior of the algorithm varies significantly depending on the value of $\eta$:
		
		\begin{enumerate}
			\item \textbf{Small Learning Rate ($0 < \eta < 0.5$)}: The term $|1 - 2\eta| < 1$ and is positive. The sequence $\{x_t\}$ converges monotonically to the minimum (0). However, if $\eta$ is very small, convergence is extremely slow.
			\item \textbf{Optimal Learning Rate ($\eta = 0.5$)}: The term $1 - 2\eta = 0$. The algorithm converges to the exact minimum in a single step ($x_1 = 0$).
			\item \textbf{Moderate Learning Rate ($0.5 < \eta < 1$)}: The term $-1 < 1 - 2\eta < 0$. The algorithm oscillates around the minimum but still converges since the magnitude of $x_t$ decreases over time.
			\item \textbf{Critical Learning Rate ($\eta = 1$)}: The term $1 - 2\eta = -1$. The sequence oscillates between $x_0$ and $-x_0$ indefinitely without ever converging.
			\item \textbf{Large Learning Rate ($\eta > 1$)}: The term $|1 - 2\eta| > 1$. The algorithm diverges as the steps become increasingly larger, moving away from the minimum.
		\end{enumerate}
		
	\end{homeworkProblem}
	
\end{document}
